\documentclass[11pt]{article}
\usepackage{mathblatt}
\newcommand{\bereich}[1]{\par\addvspace{14pt}\noindent{\large\bfseries #1}\par\nobreak\vspace{4pt}}
\newcommand{\bl}[1]{\par\noindent #1\par\vspace{3pt}}
\newcommand{\blrand}[1]{\par\noindent{\color{mbgrau}#1}\par\vspace{3pt}}
\newcommand{\blhier}[1]{\par\noindent\llap{$\blacktriangleright$\hspace{0.5em}}\textbf{#1}\par\vspace{3pt}}
\newcommand{\stufe}[1]{\par\noindent\hspace*{1.5em}{\small #1}\par\vspace{2pt}}
\newcommand{\pkt}{\textperiodcentered{}}
\begin{document}

\pfheftstil
\pflaufkopf[Abitur GK \pkt{} Prüfung]{Ableitung + Tangente}{AA30}\pfseitevon
\pfheftkopf{Ableitung + Tangente}{}
\blhier{Änderungsrate}
\stufe{mittlere Änderungsrate gesucht}
\stufe{momentane Änderungsrate gesucht}
\stufe{Stelle zur Änderungsrate gesucht}
\bl{Tangente}
\bl{Winkel am Graphen}
\blrand{Dreieck aus Tangente und Normale}

\par\addvspace{14pt}\noindent Vorher: Kurvenuntersuchung\par
\par\noindent Weiter: Integral\par
\end{document}